\begin{enumerate}
% พีชคณิตของเวกเตอร์
\item \textbf{พีชคณิตของเวกเตอร์} 

\begin{enumerate}
\item กำหนด $\vec{v} = \bv{1,2}$ และ $\vec{u} = \bv{3,4}$ จงหา
\begin{tasks}(3)
\task $2\vec{v}$
\task $\vec{v} + \vec{u}$
\task $3\vec{u} - \vec{v}$
\end{tasks}
\vspace{30pt}
\item กำหนด $\vec{v} = \bv{0,1,2}$ และ $\vec{u} = \bv{3,4,5}$ จงหา
\begin{tasks}(3)
\task $\frac{1}{2}\vec{v}$
\task $\vec{v} + \vec{u}$
\task $\frac{1}{3}\vec{v} + \frac{2}{3}\vec{u}$
\end{tasks}
\vspace{30pt}
\end{enumerate}

% เวกเตอร์ที่มีจุดเริ่มต้นไม่ใช่จุดกำเนิด
\item \textbf{เวกเตอร์ที่มีจุดเริ่มต้นไม่ใช่จุดกำเนิด} 

กำหนดจุด $P_1 = (x_1,y_1)$ และ $P_2 = (x_2, y_2)$ แล้วเวกเตอร์ที่เริ่มต้นที่จุด $P_1$ และสิ้นสุดที่จุด $P_2$ เขียนได้ในรูป

\[\rule{30pt}{.5pt} = \rule{100pt}{.5pt}\]

จงหาเวกเตอร์ที่เริ่มจากจุด $P_1$ และสิ้นสุดที่จุด $P_2$ ดังต่อไปนี้
\begin{tasks}(3)
\task $P_1 = (0,1),  P_2 = (3,2)$
\task $P_1 = (-2,1), P_2 = (1,-2)$
\task $P_1 = (5,2),  P_2 = (-1,-2)$
\end{tasks}
\vspace{50pt}

% ขนาดของเวกเตอร์
\item \textbf{ขนาดของเวกเตอร์} 

กำหนดเวกเตอร์ในปริภูมิสองมิติ $\vec{v} = \bv{v_1,v_2}$ แล้วขนาดของเวกเตอร์มีสูตรดังนี้ 

\[\rule{30pt}{.5pt} = \rule{100pt}{.5pt}\]

กำหนดเวกเตอร์ในปริภูมิสามมิติ $\vec{v} = \bv{v_1,v_2,v_3}$ แล้วขนาดของเวกเตอร์มีสูตรดังนี้

\[\rule{30pt}{.5pt} = \rule{100pt}{.5pt}\]

จงหาขนาดของเวกเตอร์ต่อไปนี้
\begin{tasks}(4)
\task $\bv{4,3}$
\task $\bv{-2, -1}$
\task $\bv{1,-2,2}$
\task $\bv{-4,0,1}$
\end{tasks}
\vspace{50pt}

\newpage
% ขนาดของเวกเตอร์
\item \textbf{เวกเตอร์หนึ่งหน่วย} 

กำหนดเวกเตอร์ในปริภูมิสองมิติ $\vec{v} = \bv{v_1,v_2}$ แล้วเวกเตอร์หนึ่งหน่วยที่มีทิศทางเดียวกันกับ $\vec{v}$ คือ \hrulefill

จงหาเวกเตอร์หนึ่งหน่วย ที่มีทิศทางเดียวกันกับเวกเตอร์ที่กำหนดให้ต่อไปนี้
\begin{tasks}(4)
\task $\bv{3,4}$
\task $\bv{0, -2}$
\task $\bv{1,2,2}$
\task $\bv{1,1,1}$
\end{tasks}
\vspace{50pt}

% ผลคูณจุด
\item \textbf{ผลคูณจุด} 

กำหนดเวกเตอร์ในปริภูมิสองมิติ $\vec{v} = \bv{v_1,v_2}$ และ $\vec{u} = \bv{u_1, u_2}$ แล้วผลคูณจุดมีสูตรดังนี้ \hrulefill

กำหนดเวกเตอร์ในปริภูมิสามมิติ $\vec{v} = \bv{v_1,v_2,v_3}$ และ $\vec{u} = \bv{u_1, u_2,u_3}$ แล้วผลคูณจุดมีสูตรดังนี้ \hrulefill

\mathnote{\textbf{สำคัญ} ผลคูณจุดเป็นปริมาณ \rule{50pt}{.5pt} ไม่ใช่ \rule{50pt}{.5pt} !}

จงคำนวณผลคูณจุดต่อไปนี้
\begin{tasks}(4)
\task $\bv{0,3} \cdot \bv{1,2}$
\task $\bv{-1, 2} \cdot \bv{3, -2}$
\task $\bv{0,1,3} \cdot \bv{2,4,0}$
\task $\bv{1,2,3} \cdot \bv{3,-3,1}$
\end{tasks}
\vspace{80pt}
\end{enumerate}
